\documentclass{amsart}

\usepackage{amsmath,amssymb,amsthm}
\usepackage[hidelinks]{hyperref}

\newtheorem{theorem}{Theorem}[section]
\newtheorem{lemma}[theorem]{Lemma}
\newtheorem{proposition}[theorem]{Proposition}
\newtheorem{corollary}[theorem]{Corollary}
\theoremstyle{definition}
\newtheorem{definition}[theorem]{Definition}
\newtheorem{question}[theorem]{Question}
\theoremstyle{remark}
\newtheorem{remark}[theorem]{Remark}

\raggedbottom

\newcommand{\C}{\mathbb{C}}
\newcommand{\Z}{\mathbb{Z}}

\begin{document}

\title[Proximity for forbidden labels in SIBO matroids]{Optimal proximity for finite forbidden labels in SIBO matroids}

\author{Knowledge Sense}
\address{Knowledge Sense}
\email{math@knowledgesense.jp}

\date{9 October 2026}

\subjclass[2020]{Primary 05B35; Secondary 20K01, 90C27}

\keywords{matroid, group-labeled matroid, proximity, forbidden labels, subsequence-interchangeable base orderability, characters of abelian groups}

\begin{abstract}
Let $M$ be a finite subsequence-interchangeably base orderable (SIBO) matroid. For $t=1,\dots,k$ let $\psi_t$ be a labeling of the ground set of $M$ by an abelian group $\Gamma_t$, and let $F_t\subseteq\Gamma_t$ be finite. Call a basis feasible if, for every $t$, the sum of its $\psi_t$-labels lies outside $F_t$. We prove that if a feasible basis exists, then every basis can be turned into a feasible basis by exchanging at most $\prod_{t=1}^k(|F_t|+1)-1$ elements, and that this bound is attained for every tuple of sizes $(|F_1|,\dots,|F_k|)$. For singleton forbidden sets the bound is $2^k-1$; this answers a question of Garamv\"olgyi, Mizutani, Oki, Schwarcz and Yamaguchi for SIBO matroids, including uniform and strongly base orderable ones, and improves their bound $\lfloor(e-1/2)k!\rfloor-1$. The proof rests on a sharp statement about consecutive sums in products of abelian groups, proved by interpolation with characters and a rank bound for a triangular matrix.
\end{abstract}

\maketitle

\section{Introduction}

Let $M$ be a matroid on a finite ground set $E$, let $\psi\colon E\to\Gamma$ be a labeling by an abelian group, and for $X\subseteq E$ write $\psi(X)$ for the sum of the labels $\psi(e)$, $e\in X$. For a finite set $F\subseteq\Gamma$ of forbidden labels, a basis $B$ is $F$-avoiding if $\psi(B)\notin F$. Such constraints express, for example, parity and congruency conditions: with $\Gamma=\Z/m\Z$ and $F=\Gamma\setminus\{c\}$, a basis is $F$-avoiding exactly when its label sum is congruent to $c$ modulo $m$. They were studied for matroid bases by Liu and Xu~\cite{LX} and by H\"orsch, Imolay, Mizutani, Oki and Schwarcz~\cite{HIMOS,HIMOS-full}. The latter authors proposed the following \emph{proximity conjecture}: if an $F$-avoiding basis exists, then every basis $A$ can be transformed into an $F$-avoiding basis $B$ with $|A\setminus B|\le|F|$ \cite[Conjecture~5.1]{HIMOS-full}. The bound $|F|$ is attained for every value of $|F|$ \cite[Section~5]{HIMOS-full}, and the conjecture implies a polynomial-time algorithm for finding an $F$-avoiding basis when $|F|$ is fixed. It is known in several cases, for example when $|F|\le2$ \cite{HIMOS-full}, when $\Gamma$ has prime order \cite{LX}, and for subsequence-interchangeably base orderable matroids, in particular strongly base orderable ones \cite[Corollary~5.15 and Remark~5.16]{HIMOS-full}, \cite[Lemma~2.4]{GMOSY}. Garamv\"olgyi, Mizutani, Oki, Schwarcz and Yamaguchi~\cite{GMOSY} proved it for sparse paving matroids and for $|F|\le4$.

Garamv\"olgyi et al.~\cite{GMOSY} also considered several constraints at once. Given labelings $\psi_t\colon E\to\Gamma_t$ and elements $f_t\in\Gamma_t$ for $t=1,\dots,k$, one asks for a basis $B$ with $\psi_t(B)\ne f_t$ for all $t$, and the multi-labeled proximity conjecture \cite[Conjecture~1.6]{GMOSY} asserts that such a basis can be reached from any basis by at most $d(k)$ exchanges, for some function $d$. They showed that necessarily $d(k)\ge2^k-1$, even for uniform matroids \cite[Example~5.1]{GMOSY}, and proved the upper bounds $3$ for $k=2$ and arbitrary matroids and $\lfloor(e-1/2)k!\rfloor-1$ for SIBO matroids \cite[Theorem~5.5]{GMOSY}. They asked whether $2^k-1$ suffices \cite[Question~6.2]{GMOSY}, adding that they did not expect this for general matroids but considered it more likely for uniform, strongly base orderable and SIBO matroids.

In this paper we settle this question for SIBO matroids, in the more general setting where each labeling has its own finite set of forbidden labels. Throughout, $\Gamma_1,\dots,\Gamma_k$ are abelian groups, $\psi_t\colon E\to\Gamma_t$ are labelings and $F_t\subseteq\Gamma_t$ are finite sets. We put $q_t=|F_t|$ and
\[
D=\prod_{t=1}^k(q_t+1),
\]
with the convention that the empty product is $1$. A set $X\subseteq E$ is \emph{feasible} if $\psi_t(X)\notin F_t$ for every $t$.

\begin{theorem}\label{thm:main}
Let $M$ be a SIBO matroid on a finite set $E$, and let $\psi_t$ and $F_t$ $(t=1,\dots,k)$ be as above. If $M$ has a feasible basis, then for every basis $A$ of $M$ there is a feasible basis $B$ with
\[
|A\setminus B|\le D-1.
\]
\end{theorem}

The definition of SIBO matroids is recalled in Section~\ref{sec:prelim}; strongly base orderable matroids, and hence all gammoids, in particular uniform, partition, laminar and transversal matroids \cite[Section~42.6c]{Schrijver}, are SIBO. Since $|A\setminus B|$ never exceeds $\min\{r,n-r\}$ for a matroid of rank $r$ on $n$ elements, the bound of Theorem~\ref{thm:main} may of course be combined with this trivial one. The groups $\Gamma_t$ are arbitrary; empty forbidden sets contribute a factor $1$ to $D$, and the case $k=0$ is included.

The bound depends only on the size tuple $(q_1,\dots,q_k)$, and among such bounds it is optimal.

\begin{theorem}\label{thm:sharp}
Let $k\ge0$ and let $q_1,\dots,q_k$ be nonnegative integers. Put $D=\prod_{t=1}^k(q_t+1)$ and $r=D-1$. There are finite cyclic groups $\Gamma_t$, labelings $\psi_t$ of the ground set of the uniform matroid $U_{r,2r}$, sets $F_t\subseteq\Gamma_t$ with $|F_t|=q_t$, and a basis $A$ of $U_{r,2r}$, such that $U_{r,2r}$ has exactly one feasible basis $B$ and $|A\setminus B|=D-1$.
\end{theorem}

In Theorem~\ref{thm:sharp} the groups and the rank $r=D-1$ are chosen according to the size tuple. For $k=1$ the example is the one given in \cite[Section~5]{HIMOS-full}, and for $q_1=\dots=q_k=1$ it is a variant of \cite[Example~5.1]{GMOSY} in which all groups are finite.

For singleton forbidden sets we obtain the following.

\begin{corollary}\label{cor:singletons}
Let $M$ be a SIBO matroid, let $\psi_t\colon E\to\Gamma_t$ be labelings and $f_t\in\Gamma_t$ for $t=1,\dots,k$. If $M$ has a basis $B'$ with $\psi_t(B')\ne f_t$ for all $t$, then for every basis $A$ there is a basis $B$ with $\psi_t(B)\ne f_t$ for all $t$ and $|A\setminus B|\le2^k-1$. The bound $2^k-1$ is attained by uniform matroids.
\end{corollary}

Thus \cite[Question~6.2]{GMOSY} has an affirmative answer for SIBO matroids and every $k$. For $k\le2$ the bound $2^k-1$ was already known to hold for all matroids \cite{HIMOS-full,GMOSY}; for larger $k$ the comparison with \cite[Theorem~5.5]{GMOSY} is shown in Table~\ref{tab:compare}. The question remains open for matroids that are not SIBO.

\begin{table}[ht]
\caption{Upper bounds on the number of exchanges for SIBO matroids and $k$ singleton forbidden sets.}\label{tab:compare}
\begin{tabular}{l|cccccc}
$k$ & 1 & 2 & 3 & 4 & 5 & 6\\
\hline
$\lfloor(e-1/2)k!\rfloor-1$ \cite{GMOSY} & 1 & 3 & 12 & 52 & 265 & 1596\\
$2^k-1$ (Corollary~\ref{cor:singletons}) & 1 & 3 & 7 & 15 & 31 & 63
\end{tabular}
\end{table}

When some $F_t$ has more than one element, Theorem~\ref{thm:main} is stronger than what can be obtained by treating each forbidden label as a separate constraint. Splitting every $F_t$ into singletons produces $\sum_t q_t$ constraints, and the bound $2^{\sum_t q_t}-1$ that \cite[Question~6.2]{GMOSY} asks for in this situation exceeds $D-1$ as soon as some $q_t\ge2$. For instance, the size tuple $(2,3)$ gives $11$ instead of $31$. Combining the labelings into one labeling with values in $\Gamma_1\times\dots\times\Gamma_k$ does not help either: the corresponding forbidden set $\{z: z_t\in F_t\text{ for some }t\}$ has a size that depends on the groups, and it is infinite as soon as $k\ge2$, some $F_t$ is nonempty and some $\Gamma_u$ with $u\ne t$ is infinite.

\subsection*{The method}
The SIBO property allows us to compare a closest feasible basis $B$ with a given basis $A$ through the bases obtained from $B$ by exchanging consecutive blocks with respect to corresponding orderings $x_1,\dots,x_m$ of $B\setminus A$ and $y_1,\dots,y_m$ of $A\setminus B$. Each of these bases is closer to $A$ and therefore infeasible. The labels enter only through the differences $d_i$ of the exchanged pairs, and the problem becomes the following statement about consecutive sums, which may be of independent interest.

\begin{theorem}\label{thm:interval}
Let $\Gamma_1,\dots,\Gamma_k$ be abelian groups, let $T_t\subseteq\Gamma_t$ be finite sets with $0\notin T_t$, and put $D=\prod_{t=1}^k(|T_t|+1)$. Let $d_1,\dots,d_\ell$ be elements of $\Gamma_1\times\dots\times\Gamma_k$, and suppose that for all $1\le i\le j\le\ell$ there is an index $t$ such that the $t$-th coordinate of $d_i+\dots+d_j$ lies in $T_t$. Then $\ell\le D-1$.
\end{theorem}

Theorem~\ref{thm:interval} is sharp for every tuple $(|T_1|,\dots,|T_k|)$; see Proposition~\ref{prop:interval-sharp}. In the language of the proof of \cite[Lemma~5.2]{GMOSY}, it shows that the constant $c_k$ defined there equals $2^k$ (Remark~\ref{rem:ck}). Our proof of Theorem~\ref{thm:interval} constructs, from complex characters, a function on a finitely generated subgroup of $\Gamma_1\times\dots\times\Gamma_k$ that equals $1$ at $0$, vanishes whenever some coordinate lies in $T_t$, and is a linear combination of at most $D$ characters. Evaluating it on the differences of the prefix sums of $d_1,\dots,d_\ell$ gives a triangular matrix of size $\ell+1$ with nonzero diagonal, whose rank is at most $D$.

\subsection*{Organization}
Section~\ref{sec:prelim} recalls SIBO matroids and reduces an SI-ordering to the symmetric difference of two bases. Section~\ref{sec:interval} proves Theorem~\ref{thm:interval} and shows that it is sharp. Section~\ref{sec:proofs} proves Theorems~\ref{thm:main} and~\ref{thm:sharp} and Corollary~\ref{cor:singletons}. Section~\ref{sec:remarks} discusses consequences for weakly base orderable matroids, an algorithmic remark and open questions.

\section{SIBO matroids}\label{sec:prelim}

We use standard matroid terminology~\cite{Oxley}. For a set $X$ and elements $x\in X$, $y\notin X$ we write $X-x=X\setminus\{x\}$ and $X+y=X\cup\{y\}$. For bases $A$ and $B$ of a matroid we have $|A\setminus B|=|B\setminus A|$, and we call this number the distance between $A$ and $B$.

\begin{definition}[\cite{Baumgart}; see {\cite[p.~85:3]{GMOSY}}]
A matroid $M$ of rank $r$ is \emph{subsequence-interchangeably base orderable} (SIBO) if for every ordered pair $(A,B)$ of bases there are orderings $a_1,\dots,a_r$ of $A$ and $b_1,\dots,b_r$ of $B$ such that
\[
(B\setminus\{b_i,\dots,b_j\})\cup\{a_i,\dots,a_j\}
\]
is a basis for all $1\le i\le j\le r$. Such a pair of orderings is an \emph{SI-ordering} of $(A,B)$.
\end{definition}

Only consecutive intervals are required. A matroid is strongly base orderable if for any two bases $A$ and $B$ there is a bijection $\varphi\colon B\to A$ such that $(B\setminus Z)\cup\varphi(Z)$ is a basis for every $Z\subseteq B$; ordering $B$ arbitrarily and $A$ through $\varphi$ shows that strongly base orderable matroids are SIBO. Baumgart~\cite{Baumgart} conjectured that graphic matroids are SIBO, and the matroid $R_{10}$ is not SIBO \cite[Theorem~4.2]{GMOSY}. Every SI-ordering yields a sequential symmetric exchange in the sense of Gabow~\cite{Gabow}.

The orderings in the definition involve all elements of $A$ and $B$, including common ones. The following lemma removes the common elements.

\begin{lemma}\label{lem:reduce}
Let $A$ and $B$ be bases of a SIBO matroid and let $m=|A\setminus B|$. There are orderings $x_1,\dots,x_m$ of $B\setminus A$ and $y_1,\dots,y_m$ of $A\setminus B$ such that
\[
B_{i,j}=(B\setminus\{x_i,\dots,x_j\})\cup\{y_i,\dots,y_j\}
\]
is a basis for all $1\le i\le j\le m$.
\end{lemma}

\begin{proof}
Let $a_1,\dots,a_r$ and $b_1,\dots,b_r$ be an SI-ordering of $(A,B)$. If $a_p\in B$ and $a_p\ne b_p$, then $a_p\in B-b_p$, so the set $(B-b_p)\cup\{a_p\}$ obtained from the interval $\{p\}$ has only $r-1$ elements and is not a basis. Hence every element of $A\cap B$ occupies the same position in both orderings, that is, $a_p\in B$ implies $a_p=b_p$. Conversely, if $b_p\in A$, then $b_p=a_{p'}$ for some $p'$; since $a_{p'}\in B$, we get $a_{p'}=b_{p'}$, hence $b_p=b_{p'}$ and $p=p'$ because $b_1,\dots,b_r$ are distinct. Thus $b_p\in A$ implies $b_p=a_p$. Let $p_1<\dots<p_m$ be the positions $p$ with $a_p\notin B$, and put $x_s=b_{p_s}$ and $y_s=a_{p_s}$. These are orderings of $B\setminus A$ and $A\setminus B$. For $1\le i\le j\le m$, the interval of positions from $p_i$ to $p_j$ contains, besides $p_i,\dots,p_j$, only positions $p$ with $a_p=b_p$, and such an element is removed from $B$ and added back. Therefore
\[
\bigl(B\setminus\{b_p: p_i\le p\le p_j\}\bigr)\cup\{a_p: p_i\le p\le p_j\}=B_{i,j},
\]
and $B_{i,j}$ is a basis.
\end{proof}

\section{Consecutive sums avoiding forbidden sets}\label{sec:interval}

A \emph{character} of an abelian group $L$ is a homomorphism $\chi\colon L\to\C^\times$; characters need not be unitary. The trivial character is the constant function $1$, and the pointwise product of two characters is a character.

\begin{lemma}\label{lem:interpolation}
Let $L$ be a finitely generated abelian group and let $T\subseteq L$ be a finite set with $0\notin T$ and $|T|=q$. There are characters $\chi_0,\dots,\chi_q$ of $L$ and complex numbers $c_0,\dots,c_q$ such that the function $\phi=\sum_{i=0}^q c_i\chi_i$ satisfies $\phi(0)=1$ and $\phi(a)=0$ for all $a\in T$.
\end{lemma}

\begin{proof}
We first show that characters separate the points of $L$, that is, for every $z\ne0$ in $L$ there is a character $\chi$ with $\chi(z)\ne1$. By the structure theorem we may assume $L=\Z^a\times\Z/n_1\Z\times\dots\times\Z/n_b\Z$. If $z$ has a nonzero coordinate $z_i$ in a factor $\Z$, the character $w\mapsto2^{w_i}$ takes the value $2^{z_i}\ne1$ at $z$. Otherwise $z$ has a coordinate $z_i\not\equiv0$ in a factor $\Z/n\Z$, and $w\mapsto\zeta^{w_i}$ with a primitive $n$-th root of unity $\zeta$ takes the value $\zeta^{z_i}\ne1$ at $z$.

Let $S=\{0\}\cup T$, so $|S|=q+1$, and let $W\subseteq\C^S$ be the linear span of the restrictions of all characters of $L$ to $S$. For distinct $a,b\in S$ choose a character $\chi_{a,b}$ with $\chi_{a,b}(a-b)\ne1$, that is, $\chi_{a,b}(a)\ne\chi_{a,b}(b)$. For $a\in S$ put
\[
h_a=\prod_{b\in S\setminus\{a\}}\frac{\chi_{a,b}-\chi_{a,b}(b)}{\chi_{a,b}(a)-\chi_{a,b}(b)}.
\]
Expanding the product writes $h_a$ as a linear combination of products of characters, that is, of characters. Its restriction to $S$ is the indicator function of $a$: every factor equals $1$ at $a$, and the factor indexed by $b$ vanishes at $b$. Hence $W=\C^S$. Choose characters $\chi_0,\dots,\chi_q$ whose restrictions to $S$ form a basis of $\C^S$, and write the indicator function of $0$ as $\sum_i c_i\chi_i|_S$. Then $\phi=\sum_i c_i\chi_i$ has the required properties.
\end{proof}

\begin{remark}
The number $q+1$ of characters in Lemma~\ref{lem:interpolation} cannot be lowered in general, and in particular a single character does not suffice. For instance, if $L=\Z/2\Z\times\Z/2\Z$ and $T=L\setminus\{0\}$, then $\phi$ must be the indicator function $\delta_0$ of $0$. The four characters of $L$ take values $\pm1$, are linearly independent, and satisfy
\[
\delta_0=\tfrac14\sum_{\chi}\chi,
\]
the sum running over all four characters. Since this representation is unique and all its coefficients are nonzero, $\delta_0$ is not a linear combination of fewer than four characters.
\end{remark}

\begin{proof}[Proof of Theorem~\ref{thm:interval}]
Write $q_t=|T_t|$ and $d_i=(d_{i,1},\dots,d_{i,k})$. For each $t$ let $L_t$ be the subgroup of $\Gamma_t$ generated by $T_t$ and $d_{1,t},\dots,d_{\ell,t}$. It is finitely generated, so Lemma~\ref{lem:interpolation} gives characters $\chi_{t,0},\dots,\chi_{t,q_t}$ of $L_t$ and coefficients $c_{t,0},\dots,c_{t,q_t}$ such that $\phi_t=\sum_i c_{t,i}\chi_{t,i}$ satisfies $\phi_t(0)=1$ and $\phi_t|_{T_t}=0$. Let $L=L_1\times\dots\times L_k$ and define $\Phi\colon L\to\C$ by
\[
\Phi(z)=\prod_{t=1}^k\phi_t(z_t).
\]
Then $\Phi(0)=1$, and $\Phi(z)=0$ whenever $z_t\in T_t$ for some $t$. Expanding the product gives
\[
\Phi=\sum_{\lambda\in\Lambda}c_\lambda\chi_\lambda,\qquad \Lambda=\prod_{t=1}^k\{0,\dots,q_t\},
\]
where $c_\lambda=\prod_t c_{t,\lambda_t}$ and $\chi_\lambda(z)=\prod_t\chi_{t,\lambda_t}(z_t)$ is a character of $L$. Note that $|\Lambda|=D$.

Put $s_0=0$ and $s_j=d_1+\dots+d_j$ for $1\le j\le\ell$; all $s_j$ lie in $L$. Consider the $(\ell+1)\times(\ell+1)$ complex matrix
\[
K=\bigl(\Phi(s_j-s_i)\bigr)_{0\le i,j\le\ell}.
\]
Its diagonal entries are $\Phi(0)=1$. If $i<j$, then $s_j-s_i=d_{i+1}+\dots+d_j$, and by hypothesis some coordinate of this element lies in the corresponding set $T_t$; hence $K_{i,j}=0$. Thus $K$ is lower triangular with unit diagonal, and $\operatorname{rank}K=\ell+1$. On the other hand, $\chi_\lambda(s_j-s_i)=\chi_\lambda(-s_i)\chi_\lambda(s_j)$, so
\[
K=\sum_{\lambda\in\Lambda}c_\lambda\,u_\lambda v_\lambda^{\mathsf T},\qquad u_\lambda=\bigl(\chi_\lambda(-s_i)\bigr)_{i=0}^{\ell},\quad v_\lambda=\bigl(\chi_\lambda(s_j)\bigr)_{j=0}^{\ell},
\]
is a sum of $D$ matrices of rank at most one. Therefore $\ell+1\le D$.
\end{proof}

The characters are used only on the finitely generated subgroups $L_t$, so no extension to $\Gamma_t$ is needed.

We now show that Theorem~\ref{thm:interval} is sharp. For nonnegative integers $q_1,\dots,q_k$ put
\begin{equation}\label{eq:N}
N_0=1,\qquad N_t=\prod_{u=1}^t(q_u+1)\quad(1\le t\le k),
\end{equation}
so that $N_k=D$, and define
\begin{equation}\label{eq:T}
T_t=\{\,jN_{t-1}+N_t\Z : 1\le j\le q_t\,\}\subseteq\Z/N_t\Z.
\end{equation}
Since $0<jN_{t-1}\le q_tN_{t-1}<N_t$ for $1\le j\le q_t$, the set $T_t$ has exactly $q_t$ elements and does not contain $0$.

\begin{lemma}\label{lem:radix}
Let $\nu$ be an integer with $1\le\nu\le D-1$. Then there is an index $t$ with $\nu+N_t\Z\in T_t$.
\end{lemma}

\begin{proof}
Since $N_0=1$ divides $\nu$ and $N_k=D$ does not, there is a smallest $t\ge1$ with $N_t\nmid\nu$. Then $N_{t-1}$ divides both $\nu$ and $N_t$, so the residue of $\nu$ modulo $N_t$ in $\{0,\dots,N_t-1\}$ equals $jN_{t-1}$ for some $0\le j<N_t/N_{t-1}=q_t+1$. As $N_t\nmid\nu$, we have $j\ge1$, and hence $\nu+N_t\Z\in T_t$.
\end{proof}

\begin{proposition}\label{prop:interval-sharp}
For all nonnegative integers $q_1,\dots,q_k$ there are finite cyclic groups $\Gamma_t$, sets $T_t\subseteq\Gamma_t$ with $0\notin T_t$ and $|T_t|=q_t$, and elements $d_1,\dots,d_{D-1}$ satisfying the hypothesis of Theorem~\ref{thm:interval}.
\end{proposition}

\begin{proof}
Take $\Gamma_t=\Z/N_t\Z$, the sets $T_t$ of \eqref{eq:T}, and $d_1=\dots=d_{D-1}=(1,\dots,1)$. For $1\le i\le j\le D-1$, the $t$-th coordinate of $d_i+\dots+d_j$ is the residue of $\nu=j-i+1$ modulo $N_t$, and $1\le\nu\le D-1$. Lemma~\ref{lem:radix} completes the proof.
\end{proof}

\begin{remark}\label{rem:ck}
In the proof of \cite[Lemma~5.2]{GMOSY}, $c_k$ denotes the least integer such that for any labelings $\psi_1,\dots,\psi_k$, elements $f_1,\dots,f_k$, set $B$ with $\psi_t(B)\ne f_t$ for all $t$, and pairwise disjoint nonempty sets $X_1,\dots,X_\ell\subseteq B$ and $Y_1,\dots,Y_\ell$ disjoint from $B$ with $\ell\ge c_k$, some interval exchange $B_{i,j}=(B\setminus(X_i\cup\dots\cup X_j))\cup(Y_i\cup\dots\cup Y_j)$ satisfies $\psi_t(B_{i,j})\ne f_t$ for all $t$; it is shown there that $c_k\le\lfloor(e-1/2)k!\rfloor$. Applying Theorem~\ref{thm:interval} with $T_t=\{f_t-\psi_t(B)\}$ and $d_i=(\psi_t(Y_i)-\psi_t(X_i))_t$ gives $c_k\le2^k$. Conversely, let $\ell=2^k-1$, let $d_1,\dots,d_\ell$ and $T_t$ be as in Proposition~\ref{prop:interval-sharp} for $q_1=\dots=q_k=1$, and let $f_t$ be the unique element of $T_t$. Take $B=\{x_1,\dots,x_\ell\}$, $X_i=\{x_i\}$ and $Y_i=\{y_i\}$, with $\psi_t(x_i)=0$ and $\psi_t(y_i)$ equal to the $t$-th coordinate of $d_i$. Then $\psi_t(B)=0\ne f_t$ for all $t$, while $\psi_t(B_{i,j})$ is the $t$-th coordinate of $d_i+\dots+d_j$, so every interval exchange has $\psi_t(B_{i,j})=f_t$ for some $t$. Hence $c_k=2^k$.
\end{remark}

\section{Proofs of the main results}\label{sec:proofs}

\begin{proof}[Proof of Theorem~\ref{thm:main}]
Fix a basis $A$. Since $M$ has a feasible basis and only finitely many bases, we may choose a feasible basis $B$ for which $m=|A\setminus B|$ is minimum. If $m=0$ there is nothing to prove, so assume $m\ge1$. Lemma~\ref{lem:reduce}, applied to $A$ and $B$, gives orderings $x_1,\dots,x_m$ of $B\setminus A$ and $y_1,\dots,y_m$ of $A\setminus B$ such that $B_{i,j}=(B\setminus\{x_i,\dots,x_j\})\cup\{y_i,\dots,y_j\}$ is a basis for all $1\le i\le j\le m$. Since $|A\setminus B_{i,j}|=m-(j-i+1)<m$, the minimality of $m$ shows that no $B_{i,j}$ is feasible.

For $1\le t\le k$ put
\[
T_t=\{\,f-\psi_t(B): f\in F_t\,\}\subseteq\Gamma_t,
\]
so that $|T_t|=q_t$ and, since $B$ is feasible, $0\notin T_t$. Let $d_i=(\psi_t(y_i)-\psi_t(x_i))_{t=1}^k\in\Gamma_1\times\dots\times\Gamma_k$. For $1\le i\le j\le m$ we have
\[
\psi_t(B_{i,j})=\psi_t(B)+\sum_{s=i}^j\bigl(\psi_t(y_s)-\psi_t(x_s)\bigr),
\]
and $B_{i,j}$ is infeasible, so $\psi_t(B_{i,j})\in F_t$ for some $t$, that is, the $t$-th coordinate of $d_i+\dots+d_j$ lies in $T_t$. Theorem~\ref{thm:interval} yields $m\le D-1$.
\end{proof}

\begin{proof}[Proof of Theorem~\ref{thm:sharp}]
Let $N_t$ and $T_t$ be as in \eqref{eq:N} and \eqref{eq:T}, so that $r=D-1=N_k-1$. Let $E$ be the disjoint union of sets $A$ and $B$ with $|A|=|B|=r$, and consider $U_{r,2r}$ on $E$; if $D=1$ this is the matroid on the empty set whose only basis is empty. For $1\le t\le k$ let $\Gamma_t=\Z/N_t\Z$, $F_t=T_t$, and
\[
\psi_t(e)=\begin{cases}1&\text{if }e\in A,\\ 0&\text{if }e\in B.\end{cases}
\]
Then $|F_t|=q_t$. For a basis $X$ put $\ell(X)=|X\cap A|$, so that $\psi_t(X)=\ell(X)+N_t\Z$ for every $t$. The only basis with $\ell(X)=0$ is $B$, and $B$ is feasible because $0\notin F_t$ for all $t$. If $1\le\ell(X)\le r=D-1$, then Lemma~\ref{lem:radix} gives an index $t$ with $\psi_t(X)\in F_t$, so $X$ is infeasible. Hence $B$ is the unique feasible basis, and $|A\setminus B|=r=D-1$.
\end{proof}

For $k=1$ this is the rank-$q_1$ uniform matroid with labels $1$ and $0$ in $\Z/(q_1+1)\Z$ and $F_1=\Z/(q_1+1)\Z\setminus\{0\}$ from \cite[Section~5]{HIMOS-full}. For $q_1=\dots=q_k=1$ we obtain $N_t=2^t$ and $F_t=\{2^{t-1}\}\subseteq\Z/2^t\Z$; the mechanism is the one of \cite[Example~5.1]{GMOSY}, where one of the groups is $\Z$.

\begin{proof}[Proof of Corollary~\ref{cor:singletons}]
Apply Theorem~\ref{thm:main} with $F_t=\{f_t\}$, so that $D=2^k$, and Theorem~\ref{thm:sharp} with $q_1=\dots=q_k=1$.
\end{proof}

\section{Further remarks}\label{sec:remarks}

\subsection{Weakly base orderable matroids}
The proof of Theorem~\ref{thm:main} uses the SIBO property only to produce many infeasible interval exchanges. Following \cite{HIMOS-full}, an ordered pair $(B_1,B_2)$ of bases has the \emph{$p$-exchange property} if there are pairwise disjoint nonempty sets $X_1,\dots,X_p\subseteq B_1\setminus B_2$ and $Y_1,\dots,Y_p\subseteq B_2\setminus B_1$ such that $(B_1\setminus\bigcup_{i\in Z}X_i)\cup\bigcup_{i\in Z}Y_i$ is a basis for every $Z\subseteq\{1,\dots,p\}$, and a matroid is \emph{weakly $(\alpha,p)$-base orderable} if every ordered pair $(B_1,B_2)$ of bases with $|B_1\setminus B_2|\ge\alpha$ has the $p$-exchange property.

\begin{corollary}\label{cor:weak}
Let $M$ be a weakly $(\alpha,D)$-base orderable matroid, and let $\psi_t$ and $F_t$ be as in Theorem~\ref{thm:main}. If $M$ has a feasible basis, then for every basis $A$ there is a feasible basis $B$ with $|A\setminus B|\le\alpha-1$.
\end{corollary}

\begin{proof}
Choose a feasible basis $B$ with $m=|A\setminus B|$ minimum, and suppose that $m\ge\alpha$. Let $X_1,\dots,X_D$ and $Y_1,\dots,Y_D$ be given by the $D$-exchange property of $(B,A)$. Taking $Z=\{i\}$ shows $|X_i|=|Y_i|$. For $1\le i\le j\le D$ the basis $(B\setminus(X_i\cup\dots\cup X_j))\cup(Y_i\cup\dots\cup Y_j)$ is at distance $m-|Y_i\cup\dots\cup Y_j|<m$ from $A$, hence infeasible. With $T_t$ as in the proof of Theorem~\ref{thm:main} and $d_i=(\psi_t(Y_i)-\psi_t(X_i))_t$, Theorem~\ref{thm:interval} gives $D\le D-1$, a contradiction.
\end{proof}

By \cite[Theorem~5.6]{GMOSY} there is a computable function $h$ such that every $\mathrm{GF}(q)$-representable matroid is weakly $(h(q,p),p)$-base orderable, and by \cite[Remark~5.10]{GMOSY} every sparse paving matroid is elementarily, and hence weakly, $(2p,p)$-base orderable. Corollary~\ref{cor:weak} therefore gives the bounds $h(q,D)-1$ and $2D-1$ for these classes. For singleton forbidden sets this replaces the number $\lfloor(e-1/2)k!\rfloor$ of blocks used in \cite[Corollary~5.11]{GMOSY} by $2^k$.

\subsection{An algorithmic remark}
Let $M$ be a SIBO matroid of rank $r$ on $n$ elements, given by an independence oracle, and assume that the group operations and equality tests in the groups $\Gamma_t$ can be performed by oracles. Given a basis $A$, one may test all sets $(A\setminus X)\cup Y$ with $X\subseteq A$, $Y\subseteq E\setminus A$ and $|X|=|Y|\le\min\{D-1,r,n-r\}$; by Theorem~\ref{thm:main} a feasible basis is found if one exists. The number of candidates is $\sum_{j=0}^{\min\{D-1,r,n-r\}}\binom{r}{j}\binom{n-r}{j}$, which is polynomial in $n$ for fixed $D$, in particular for fixed $k$ when all forbidden sets are singletons. This direct enumeration is an XP algorithm with parameter $D$.

\subsection{Questions}
For $k=1$, Theorem~\ref{thm:main} is the SIBO case of the proximity conjecture of H\"orsch et al., which was already known; our proof differs from the pigeonhole argument of \cite[Remark~5.16]{HIMOS-full} and \cite[Lemma~2.4]{GMOSY}. The following question contains both the proximity conjecture ($k=1$) and \cite[Question~6.2]{GMOSY} ($q_1=\dots=q_k=1$).

\begin{question}
Does the conclusion of Theorem~\ref{thm:main} hold for every matroid?
\end{question}

By \cite[Theorem~4.4 and Theorem~5.5]{GMOSY} the answer is affirmative when $k=1$ and $q_1\le4$, and when $k=2$ and $q_1=q_2=1$. If graphic matroids are SIBO, as conjectured by Baumgart~\cite{Baumgart}, then Theorem~\ref{thm:main} answers the question affirmatively for graphic matroids.

\section*{Use of AI-assisted tools}
AI-assisted tools were used in the research and preparation of this manuscript.

\begin{thebibliography}{9}

\bibitem{Baumgart}
M.~Baumgart, \emph{Ranking and ordering problems of spanning trees}, Ph.D. thesis, Technische Universit\"at M\"unchen, 2009.

\bibitem{Gabow}
H.~Gabow, \emph{Decomposing symmetric exchanges in matroid bases}, Math. Programming \textbf{10} (1976), no.~1, 271--276.

\bibitem{GMOSY}
D.~Garamv\"olgyi, R.~Mizutani, T.~Oki, T.~Schwarcz, and Y.~Yamaguchi, \emph{Towards the proximity conjecture on group-labeled matroids}, 52nd International Colloquium on Automata, Languages, and Programming (ICALP 2025), LIPIcs, vol.~334, Schloss Dagstuhl, Leibniz-Zentrum f\"ur Informatik, 2025, Art. No.~85, 17~pp. Full version: arXiv:2411.06771.

\bibitem{HIMOS}
F.~H\"orsch, A.~Imolay, R.~Mizutani, T.~Oki, and T.~Schwarcz, \emph{Problems on group-labeled matroid bases}, 51st International Colloquium on Automata, Languages, and Programming (ICALP 2024), LIPIcs, vol.~297, Schloss Dagstuhl, Leibniz-Zentrum f\"ur Informatik, 2024, Art. No.~86, 20~pp.

\bibitem{HIMOS-full}
F.~H\"orsch, A.~Imolay, R.~Mizutani, T.~Oki, and T.~Schwarcz, \emph{Problems on group-labeled matroid bases}, preprint (2024), arXiv:2402.16259v1. Full version of~\cite{HIMOS}.

\bibitem{LX}
S.~Liu and C.~Xu, \emph{On the congruency-constrained matroid base}, Integer Programming and Combinatorial Optimization (IPCO 2024), Lecture Notes in Comput. Sci., vol.~14679, Springer, Cham, 2024, pp.~280--293, doi:10.1007/978-3-031-59835-7\_21.

\bibitem{Oxley}
J.~Oxley, \emph{Matroid theory}, 2nd ed., Oxford Graduate Texts in Mathematics, vol.~21, Oxford University Press, Oxford, 2011.

\bibitem{Schrijver}
A.~Schrijver, \emph{Combinatorial optimization: polyhedra and efficiency}, Springer, Berlin, 2003.

\end{thebibliography}

\end{document}
